\documentclass{beamer}

% --- Tema: simple y limpio, sin colores estridentes ---
\usetheme{Madrid}
\setbeamertemplate{navigation symbols}{}

% --- Codificación ---
\usepackage[utf8]{inputenc}
\usepackage[T1]{fontenc}
\usepackage[spanish]{babel}

% --- Diagramas ---
\usepackage{tikz}
\usepackage{graphicx}
\usetikzlibrary{arrows.meta, positioning}

% --- Fragmentos de código (misma config que informe/informe.tex) ---
\usepackage{listings}
\usepackage{xcolor}
\lstset{
    basicstyle=\ttfamily\small,
    breaklines=true,
    frame=single,
    numbers=left,
    numbersep=5pt,
    commentstyle=\color{gray},
    keywordstyle=\color{blue},
    stringstyle=\color{teal},
    extendedchars=true,
    literate=
     {á}{{\'a}}1 {é}{{\'e}}1 {í}{{\'i}}1 {ó}{{\'o}}1 {ú}{{\'u}}1
     {Á}{{\'A}}1 {É}{{\'E}}1 {Í}{{\'I}}1 {Ó}{{\'O}}1 {Ú}{{\'U}}1
     {ñ}{{\~n}}1 {Ñ}{{\~N}}1 {¿}{{?`}}1 {¡}{{!`}}1
}

% --- Bibliografía (natbib + bibtex, reutiliza informe/referencias.bib) ---
\usepackage{natbib}

% --- Portada ---
\title[Consumo REST -- ArtiSync]{Consumo de Web Services REST desde un Front-end AlpineJS --- Proyecto ArtiSync}
\subtitle{Proyecto ArtiSync}
\author[Figueroa \& Loor]{Figueroa Morales Bryan Javier \and Loor Medranda Marlon Taylor}
\institute[UTEQ]{
    Universidad Técnica Estatal de Quevedo \\
    Facultad de Ciencias de la Computación y Diseño Digital \\
    Carrera de Ingeniería de Software \\[6pt]
    Aplicaciones Web (5to nivel) --- Segundo Corte \\
    Docente: Dr. Gleiston Cicerón Guerrero Ulloa, Ph.D.
}
\date{Julio 2026}

\begin{document}

% ============================================================
% 1. PORTADA
% ============================================================
\begin{frame}
    \titlepage
\end{frame}

% ============================================================
% 2. CONTEXTO
% ============================================================
\begin{frame}{Contexto}
    \vfill
    \begin{itemize}
        \item \textbf{ArtiSync}: plataforma que conecta clientes y creadores de servicios digitales.
        \item Primer Corte: backend Spring Boot + PostgreSQL, con su API REST.
        \item Este corte agrega un \textbf{frontend AlpineJS} sobre esa API REST \cite{fielding2000}.
        \item Alcance: autenticación (login + 2FA) y CRUD de usuarios.
    \end{itemize}
    \vfill
\end{frame}

% ============================================================
% 3. DIAGRAMA DE FLUJO
% ============================================================
\begin{frame}{Flujo general}
    \vfill
    \centering
    \resizebox{0.98\linewidth}{!}{%
    \begin{tikzpicture}[
        node distance=1.8cm and 2.2cm,
        box/.style={draw, thick, rounded corners, minimum width=3.6cm, minimum height=1.6cm, align=center, font=\normalsize}
    ]
        \node[box] (cliente) {Cliente\\AlpineJS};
        \node[box, right=of cliente] (backend) {Backend\\Spring Boot (REST)};
        \node[box, right=of backend] (bd) {Base de\\Datos};

        \draw[-{Stealth}, thick] ([yshift=4pt]cliente.east) -- node[above, font=\small]{JSON / HTTP} ([yshift=4pt]backend.west);
        \draw[-{Stealth}, thick] ([yshift=-4pt]backend.west) -- node[below, font=\small]{respuesta} ([yshift=-4pt]cliente.east);

        \draw[-{Stealth}, thick] ([yshift=4pt]backend.east) -- ([yshift=4pt]bd.west);
        \draw[-{Stealth}, thick] ([yshift=-4pt]bd.west) -- ([yshift=-4pt]backend.east);
    \end{tikzpicture}%
    }

    \vspace{16pt}
    \begin{itemize}
        \item Cliente envía la petición HTTP con JSON.
        \item Backend valida el JWT y procesa la lógica de negocio.
        \item Backend persiste o consulta datos en PostgreSQL.
    \end{itemize}
    \vfill
\end{frame}

% ============================================================
% 4. AUTENTICACIÓN
% ============================================================
\begin{frame}[fragile]{Autenticación}
    \begin{itemize}
        \item \texttt{POST /auth/login} envía correo y contraseña.
        \item La respuesta trae el campo \texttt{requiere2fa}.
        \item Si es \texttt{true}, se pide un código de 6 dígitos vía \texttt{POST /auth/2fa/verify}.
    \end{itemize}
    \vspace{4pt}
    \lstinputlisting[firstline=291,lastline=294,firstnumber=291]{../login.html}
\end{frame}

% ============================================================
% 5. CRUD DE USUARIOS
% ============================================================
\begin{frame}{CRUD de usuarios}
    \vfill
    \begin{itemize}
        \item Cuatro verbos HTTP \cite{rfc9110} sobre el recurso \texttt{usuarios}.
    \end{itemize}
    \vspace{12pt}
    \begin{center}
    \renewcommand{\arraystretch}{1.6}
    \large
    \begin{tabular}{ll}
        \hline
        \textbf{Verbo} & \textbf{Endpoint} \\
        \hline
        GET & \texttt{/admin/usuarios} \\
        PUT & \texttt{/admin/usuarios/\{id\}} \\
        PATCH & \texttt{/admin/usuarios/\{id\}/estado} \\
        DELETE & \texttt{/admin/usuarios/\{id\}} (+ \texttt{/sesiones}) \\
        \hline
    \end{tabular}
    \end{center}
    \vspace{12pt}
    \begin{itemize}
        \item Todas pasan por \texttt{fetchWithAuth()}, que agrega el token automáticamente.
    \end{itemize}
    \vfill
\end{frame}

% ============================================================
% 6. MANEJO DE ERRORES
% ============================================================
\begin{frame}{Manejo de errores}
    \vfill
    \begin{itemize}
        \item \texttt{fetchWithAuth()} centraliza la detección de sesiones inválidas.
        \item Respuesta \texttt{401} \(\Rightarrow\) se limpia el token de \texttt{localStorage}.
        \item Redirección automática a \texttt{login.html} (salvo que ya esté ahí).
        \item A nivel de formulario, los mensajes de error se toman de \texttt{data.message} enviado por el propio backend.
    \end{itemize}
    \vfill
\end{frame}

% ============================================================
% 7. CIERRE
% ============================================================
\begin{frame}{Cierre}
    \begin{itemize}
        \item Login con 2FA y CRUD completo, funcionando a nivel de frontend.
        \item Prueba end-to-end real (Docker): registro y login exitosos; CRUD admin con 403 esperado (falta rol ADMIN real).
        \item Documentación reproducible: cada paso quedó respaldado con evidencia real contra el backend en ejecución.\footnote{\scriptsize Diseño de diapositivas inspirado en \cite{alley2013} y \cite{mayer2009}.}
    \end{itemize}
    \vspace{10pt}
    \begin{center}
        \Large ¡Gracias! ¿Preguntas?
    \end{center}
\end{frame}

% ============================================================
% REFERENCIAS
% ============================================================
\begin{frame}[allowframebreaks]{Referencias}
    \footnotesize
    \bibliographystyle{plain}
    \bibliography{../informe/referencias}
\end{frame}

\end{document}
